% Generated draft: requires mathematical and source review.
\documentclass[11pt]{article}
\usepackage[margin=25mm]{geometry}
\usepackage{fontspec}
\setmainfont{Latin Modern Roman}
\usepackage{amsmath,amssymb,mathtools,unicode-math}
\setmathfont{Latin Modern Math}
\usepackage{enumitem,needspace,xurl,hyperref}
\setlength{\parindent}{0pt}
\setlength{\parskip}{5pt}
\setlength{\emergencystretch}{3em}
\newcommand{\N}{\mathbb{N}}
\newcommand{\Z}{\mathbb{Z}}
\newcommand{\Q}{\mathbb{Q}}
\newcommand{\R}{\mathbb{R}}
\newcommand{\C}{\mathbb{C}}
\newcommand{\F}{\mathbb{F}}
\newcommand{\A}{\mathbb{A}}
\newcommand{\PP}{\mathbb{P}}
\newcommand{\E}{\mathbb{E}}
\newcommand{\eps}{\varepsilon}
\newcommand{\dd}{\,\mathrm{d}}
\begin{document}
\section*{UnsolvedMath ID: 30000989}
\textbf{Source problem code:} catalogue entry
\subsection{Definitions and mathematical statement}
For a space \(X\) and field \(k\), the zeroth singular Betti number is \(b_0(X;k)=\dim_k H_0(X;k)\), with the dimension understood as a cardinal when necessary. Singular \(H_0\) counts path components; for spaces such as finite simplicial complexes or locally path-connected spaces this agrees with the number of connected components. The source uses the latter description without restricting the space class.

The retained report question asks why this topological component count is regarded as stable and easy to recover, whereas the number of clusters of a finite sampled statistical space is regarded as unstable and difficult to define or recover. This is a conceptual comparison, not a specified theorem: no definition of a statistical cluster, sampling model, input representation, perturbation metric, recovery criterion or computational resource bound is given. In particular, homotopy invariance of \(b_0\) does not itself specify stability under arbitrary data perturbations. A precise comparison would require choosing all these structures, which this extraction does not do.
\subsection{Short English statement}
The source asks conceptually why topological components seem easier to define and recover than statistical clusters. It supplies no common stability or computation model, so the draft records the question and its missing definitions rather than inventing a clustering theorem.
\subsection{Sources}
\begin{itemize}
\item Local UnsolvedMath registry, permanent ID \(30000989\), source code \texttt{catalogue entry}, ``Stability of Topological Components versus Statistical Clusters.'' Registry curation: Ulam AI / UnsolvedMath Contributors, CC BY 4.0; underlying sources retain their supplied provenance.
\item \emph{Computational Algebraic Topology} (2009), Oberwolfach Reports, pp.\ 1603--1654, DOI \url{https://doi.org/10.4171/owr/2008/29}.
\item Supplied literature triage, \texttt{literature\_checked\_at}: 2026-08-21: labeled open, but assessed as a rhetorical title-level question rather than a mathematically specified claim, with no unique proposition to verify. This is dated registry curation, requiring independent review rather than certifying a mathematical open problem.
\end{itemize}
\end{document}
